%% GENERATED by tools/gen_latex.py -- edit content/ch04.py instead.
%% Build:  latexmk -pdf ch04.tex     (or: pdflatex ch04.tex, twice)
\documentclass[aspectratio=169,11pt,t]{beamer}
\usepackage{itbpro}

\coursetitle{Klasifikasi dan Regresi}
\coursesubtitle{Tiga alur kerja lengkap -- biner, multikelas, dan regresi skalar -- beserta aturan pilih loss dan aktivasi yang dipakai sisa buku ini.}
\coursekicker{Chapter 4}
\coursesource{Chollet \& Watson, Deep Learning with Python 3e -- bab 4}
\courseduration{3 jam (2 sesi)}
\coursepresenter{Rahman Indra Kesuma, S.Kom., M.Cs.}{Asisten Pengajar}
\coursebrand{AI for Professional \textperiodcentered\ ITB \textperiodcentered\ Ch 4}

\title{Klasifikasi dan Regresi}
\author{Rahman Indra Kesuma, S.Kom., M.Cs.}
\date{}

\begin{document}
\covertitle

\framekicker{Learning outcomes}
\begin{frame}[shrink=25]{Tujuan Sesi}
\leadpar{Setelah sesi ini, peserta dapat:}
\begin{isteps}
\item Memilih \textbf{aktivasi keluaran dan fungsi rugi} yang benar untuk klasifikasi biner, multikelas, dan regresi -- tanpa menebak.
\item Menyiapkan data teks dengan \textbf{multi-hot encoding}, dan label dengan \textbf{one-hot} atau \textbf{sparse}.
\item Membaca \textbf{kurva rugi validasi} untuk menemukan titik berhenti terbaik, lalu melatih ulang sampai epoch itu saja.
\item Mengenali \textbf{leher botol informasi} dan menaksir ukuran lapis antara dari banyaknya kelas.
\item Menormalkan fitur dengan \textbf{statistik data latih}, dan menilai model berdata sedikit dengan \textbf{validasi silang K-lipat}.
\end{isteps}

\vskip 6pt
\vskip 4pt
\reslink{Buku}{Bab 4 --- teks penuh}{https://deeplearningwithpython.io/chapters/chapter04\_classification-and-regression}
\reslink{Notebook}{01\_imdb\_klasifikasi\_biner.ipynb}{../../notebooks/ch04/01\_imdb\_klasifikasi\_biner.ipynb}
\reslink{Notebook}{02\_reuters\_multikelas.ipynb}{../../notebooks/ch04/02\_reuters\_multikelas.ipynb}
\reslink{Notebook}{03\_regresi\_harga\_rumah.ipynb}{../../notebooks/ch04/03\_regresi\_harga\_rumah.ipynb}
\reslink{Repo}{Notebook resmi penulis}{https://github.com/fchollet/deep-learning-with-python-notebooks/blob/master/chapter04\_classification-and-regression.ipynb}
\vskip 2pt
\end{frame}

\framekicker{Kosakata}
\begin{frame}[shrink=25]{Dua belas istilah yang dipakai sampai bab 20}
\begin{minipage}[t]{0.485\textwidth}
\begin{itable}{>{\raggedright\arraybackslash}p{0.282\textwidth}>{\raggedright\arraybackslash}p{0.658\textwidth}}
\thead{Istilah} & \thead{Artinya} \\ \midrule
\textbf{Sample / input} & Satu titik data yang masuk ke model. \\
\textbf{Prediction / output} & Yang keluar dari model. \\
\textbf{Target} & Jawaban benar, dari sumber di luar model. \\
\textbf{Loss value} & Ukuran jarak prediksi ke target. \\
\textbf{Classes} & Himpunan label yang mungkin. \\
\textbf{Mini-batch} & Lazimnya 8-128 sampel sekaligus. \\
\end{itable}

\end{minipage}%
\hfill
\begin{minipage}[t]{0.485\textwidth}
\begin{itable}{>{\raggedright\arraybackslash}p{0.3572\textwidth}>{\raggedright\arraybackslash}p{0.5828\textwidth}}
\thead{Jenis tugas} & \thead{Cirinya} \\ \midrule
\textbf{Binary classification} & Dua kategori, saling meniadakan. \\
\textbf{Multiclass} & Lebih dari dua kategori, satu label per sampel. \\
\textbf{Multilabel} & Satu sampel boleh punya beberapa label. \\
\textbf{Scalar regression} & Satu nilai kontinu. \\
\textbf{Vector regression} & Beberapa nilai kontinu sekaligus. \\
\end{itable}

\end{minipage}%
\begin{iband}[signal]
Membedakan \textbf{multiclass} dari \textbf{multilabel} menentukan pilihan aktivasi keluaran: \hilite{softmax} untuk yang pertama (jumlahnya 1), \hilite{sigmoid per kelas} untuk yang kedua.
\end{iband}
\end{frame}

\coursesection{01}{Klasifikasi biner: ulasan IMDB}{50.000 ulasan film, positif atau negatif.}

\framekicker{Bagian 4.1}
\begin{frame}[shrink=25]{Data, dan cara mengubah deret kata jadi tensor}
\icode{python}{listing 4.1-4.3 --- muat dan vektorkan}{listings/ch04-01.py}
\ioutput{listings/ch04-02.txt}
\begin{ibullets}
\item 25.000 latih + 25.000 uji, seimbang 50:50 positif-negatif.
\item \inlinecode{num\_words=10000} hanya menyimpan \textbf{10.000 kata tersering}; sisanya dibuang.
\item Multi-hot \hilite{membuang urutan kata}. Cukup untuk sentimen; bab 14-15 menggantinya saat urutan mulai penting.
\end{ibullets}
\end{frame}

\framekicker{Bagian 4.1}
\begin{frame}[shrink=25]{Model, validasi, dan titik balik di epoch 4}
\begin{minipage}[t]{0.485\textwidth}
\icode{python}{listing 4.4-4.6}{listings/ch04-03.py}

\end{minipage}%
\hfill
\begin{minipage}[t]{0.485\textwidth}
\begin{center}
\adjustbox{max width=\linewidth,max totalheight=0.52\textheight}{%

\begin{tikzpicture}[font=\sffamily\tiny]
  \draw[rule, line width=0.8pt] (0,0) -- (8.6,0);
  \draw[rule, line width=0.8pt] (0,0) -- (0,3.0);
  \node[text=ink3, anchor=north] at (4.3,-0.15) {epoch};
  \node[text=ink3, rotate=90, anchor=south] at (-0.35,1.5) {rugi};
  \draw[signal, line width=1.2pt]
    (0,2.6) .. controls (1.2,1.7) and (2.4,1.05) .. (4.0,0.62)
            .. controls (5.6,0.35) and (7.0,0.22) .. (8.6,0.15);
  \draw[rose, line width=1.2pt]
    (0,2.4) .. controls (0.8,1.6) and (1.5,1.15) .. (2.1,1.02);
  \draw[rose, line width=1.2pt]
    (2.1,1.02) .. controls (3.6,1.25) and (5.8,1.95) .. (8.6,2.75);
  \draw[amberbr, line width=0.9pt, dashed] (2.1,0) -- (2.1,3.0);
  \fill[amberbr] (2.1,1.02) circle (2.2pt);
  \node[text=amber, anchor=west] at (2.25,2.85) {epoch 4 --- titik terbaik};
  \node[text=rose, anchor=west] at (4.2,1.95) {mulai overfitting};
  \draw[signal, line width=1.2pt] (4.6,-0.55) -- (4.95,-0.55);
  \node[text=ink3, anchor=west] at (5.0,-0.55) {rugi latih};
  \draw[rose, line width=1.2pt] (6.5,-0.55) -- (6.85,-0.55);
  \node[text=ink3, anchor=west] at (6.9,-0.55) {rugi validasi};
\end{tikzpicture}

}
\end{center}
\vskip 3pt{\color{ink3}\fontsize{7.4}{9.4}\selectfont Rugi latih turun terus; rugi validasi berbalik naik setelah epoch 4.\par}

\end{minipage}%
\begin{iquote}
\quotetext{Setelah epoch keempat, Anda mengoptimalkan berlebihan pada data latih, dan berakhir mempelajari representasi yang khas untuk data latih itu saja -- yang tidak berlaku di luarnya.}\quotecite{Chollet \& Watson, bab 4}
\end{iquote}
\end{frame}
\note{Ini pola yang akan muncul di setiap bab sesudahnya. Latih 20 epoch untuk MELIHAT titik baliknya, lalu latih ulang sampai titik itu.}

\framekicker{Bagian 4.1}
\begin{frame}[shrink=25]{Hasilnya, dan tolok banding yang harus dikalahkan}
\begin{minipage}[t]{0.3253\textwidth}
\istat{88\%}{akurasi uji setelah dilatih ulang 4 epoch}
\end{minipage}%
\hfill
\begin{minipage}[t]{0.3253\textwidth}
\istat{50\%}{tolok banding acak (kelasnya seimbang)}
\end{minipage}%
\hfill
\begin{minipage}[t]{0.3253\textwidth}
\istat{16}{unit per lapis antara -- model sengaja kecil}
\end{minipage}%
\begin{iquote}
\quotetext{Tanpa fungsi aktivasi seperti \inlinecode{relu} (disebut juga nonlinearitas), lapis \inlinecode{Dense} hanya terdiri dari dua operasi linear -- hasil kali titik dan penjumlahan. Ruang hipotesis seperti itu terlalu sempit.}\quotecite{Chollet \& Watson, bab 4}
\end{iquote}
\begin{iband}[signal]
\textbf{Crossentropy} datang dari teori informasi: ia mengukur \hilite{jarak antara dua sebaran peluang}. Itulah sebabnya ia pasangan alami bagi keluaran sigmoid dan softmax, dan bukan MSE.
\end{iband}
\end{frame}

\coursesection{02}{Multikelas: kawat berita Reuters}{46 topik, saling meniadakan.}

\framekicker{Bagian 4.2}
\begin{frame}[shrink=25]{One-hot atau sparse -- antarmuka beda, hitungannya sama}
\begin{minipage}[t]{0.485\textwidth}
\icode{python}{label one-hot}{listings/ch04-04.py}

\end{minipage}%
\hfill
\begin{minipage}[t]{0.485\textwidth}
\icode{python}{label sparse (bilangan bulat)}{listings/ch04-05.py}

\end{minipage}%
\icode{python}{listing 4.11 --- model dan metrik top-K}{listings/ch04-06.py}
\begin{iband}[signal]
\textbf{Top-K accuracy} menanyakan hal yang berbeda: apakah kelas yang benar ada di antara k tebakan teratas? Untuk 46 kelas, itu ukuran yang \hilite{jauh lebih berguna} bagi sistem yang hanya menyarankan.
\end{iband}
\end{frame}

\framekicker{Bagian 4.2}
\begin{frame}[shrink=25]{Leher botol informasi: percobaan yang sengaja digagalkan}
\begin{center}
\adjustbox{max width=\linewidth,max totalheight=0.52\textheight}{%

\begin{tikzpicture}[font=\sffamily\tiny,
  u/.style={draw=signal!60, fill=signal!9, rounded corners=3pt, minimum width=1.1cm,
            minimum height=0.75cm, text=ink},
  ar/.style={-{Stealth[length=4pt]}, signal, line width=0.7pt}]
  \node[text=lime, anchor=west] at (0,1.15) {yang benar};
  \node[u] (a1) at (0.7,0.55) {\ttfamily 64};
  \node[u] (a2) at (2.3,0.55) {\ttfamily 64};
  \node[draw=lime!60, fill=limebr!14, rounded corners=3pt, minimum width=1.2cm,
        minimum height=0.75cm, text=ink] (a3) at (3.9,0.55) {\ttfamily 46};
  \draw[ar] (a1) -- (a2); \draw[ar] (a2) -- (a3);
  \node[text=lime, anchor=west] at (4.7,0.55) {akurasi ${\sim}80\%$};

  \draw[rule] (0,0.0) -- (9.0,0.0);

  \node[text=rose, anchor=west] at (0,-0.4) {leher botol informasi};
  \node[u] (b1) at (0.7,-1.05) {\ttfamily 64};
  \node[draw=rose!75, fill=rosebr!18, rounded corners=3pt, minimum width=1.1cm,
        minimum height=0.32cm, text=ink] (b2) at (2.3,-1.05) {\ttfamily 4};
  \node[u] (b3) at (3.9,-1.05) {\ttfamily 46};
  \draw[ar] (b1) -- (b2); \draw[ar] (b2) -- (b3);
  \node[text=rose, anchor=west] at (4.7,-1.05) {akurasi ${\sim}71\%$ --- turun 8 poin};
  \node[text=ink3, anchor=west] at (4.7,-1.4) {46 kelas tidak muat diperas ke 4 dimensi};
\end{tikzpicture}

}
\end{center}
\vskip 3pt{\color{ink3}\fontsize{7.4}{9.4}\selectfont Lapis 4 unit di tengah memaksa 46 kelas diperas ke 4 dimensi. Turun 8 poin akurasi.\par}
\begin{iquote}
\quotetext{Model itu sanggup menjejalkan sebagian besar informasi yang diperlukan ke dalam representasi 4 dimensi -- tetapi tidak semuanya.}\quotecite{Chollet \& Watson, bab 4}
\end{iquote}
\begin{iband}[amber]
Aturan praktisnya: \textbf{lapis antara tidak boleh lebih sempit dari jumlah kelas keluaran}. Itulah sebabnya contoh IMDB cukup 16 unit (2 kelas), sementara Reuters butuh 64.
\end{iband}
\end{frame}
\note{Percobaan gagal ini justru bagian terbaik dari bab 4 --- tunjukkan utuh, jangan lewati. Peserta belajar lebih banyak dari yang gagal.}

\framekicker{Bagian 4.2}
\begin{frame}[shrink=25]{Melatih ulang di epoch terbaik}
\icode{python}{listing 4.14 --- model produksi}{listings/ch04-07.py}
\ioutput{listings/ch04-08.txt}
\begin{minipage}[t]{0.494\textwidth}
\istat{\textasciitilde{}80\%}{akurasi uji}
\end{minipage}%
\hfill
\begin{minipage}[t]{0.494\textwidth}
\istat{\textasciitilde{}19\%}{tolok banding acak -- 46 kelas, sebaran tak rata}
\end{minipage}%
\begin{iband}[signal]
Selalu hitung \textbf{tolok banding akal sehat} dulu. Angka 80\% terdengar biasa saja sampai Anda tahu bahwa menebak asal hanya memberi 19\%.
\end{iband}
\end{frame}

\coursesection{03}{Regresi skalar: harga rumah}{Data sedikit. Aturan mainnya berubah.}

\framekicker{Bagian 4.3}
\begin{frame}[shrink=25]{Normalisasi fitur -- dan jebakan kebocoran data}
\icode{python}{listing 4.16-4.17 --- muat dan normalkan}{listings/ch04-09.py}
\begin{iband}[rose]
Baris \inlinecode{x\_test = (test\_data - mean) / std} memakai \inlinecode{mean} dan \inlinecode{std} dari data \textbf{latih}. Menghitung ulang dari data uji adalah \hilite{kebocoran informasi}: model jadi tahu sesuatu tentang data yang seharusnya belum pernah dilihatnya. Bab 5 menamainya secara resmi.
\end{iband}
\begin{ibullets}
\item 8 fitur: bujur, lintang, umur rumah, populasi, jumlah rumah tangga, penghasilan median, total kamar, total kamar tidur.
\item Target: harga rumah median, kontinu, kira-kira \$60 ribu sampai \$500 ribu.
\item Rentang tiap fitur berbeda jauh; tanpa normalisasi \hilite{optimalisasi jadi kacau}.
\end{ibullets}
\end{frame}

\framekicker{Bagian 4.3}
\begin{frame}[shrink=25]{Model regresi: keluaran tanpa aktivasi}
\icode{python}{listing 4.18}{listings/ch04-10.py}
\begin{minipage}[t]{0.3253\textwidth}
\begin{icardaccent}
\cardh{Tanpa aktivasi keluaran}
Sigmoid akan mengurung tebakan ke [0,1]. Regresi harus bebas.
\end{icardaccent}
\end{minipage}%
\hfill
\begin{minipage}[t]{0.3253\textwidth}
\begin{icardwarn}
\cardh{Model sengaja kecil}
480 sampel saja. Model besar akan langsung menghafalnya.
\end{icardwarn}
\end{minipage}%
\hfill
\begin{minipage}[t]{0.3253\textwidth}
\begin{icardaccent}
\cardh{MSE untuk rugi, MAE untuk metrik}
MAE bisa dibaca manusia: MAE 0,5 $\approx$ meleset \$50 ribu.
\end{icardaccent}
\end{minipage}%
\end{frame}

\framekicker{Bagian 4.3}
\begin{frame}[shrink=25]{Validasi silang K-lipat, karena datanya sedikit}
\begin{center}
\adjustbox{max width=\linewidth,max totalheight=0.52\textheight}{%

\begin{tikzpicture}[font=\sffamily\tiny,
  tr/.style={draw=rule, fill=papertint, rounded corners=2pt,
             minimum width=1.55cm, minimum height=0.34cm},
  va/.style={draw=signal!70, fill=signal!22, rounded corners=2pt,
             minimum width=1.55cm, minimum height=0.34cm}]
  \node[text=ink3, anchor=west] at (0,1.65) {data latih dibagi 4; tiap lipatan sekali jadi validasi};
  \foreach \r/\v in {0/1, 1/2, 2/3, 3/4} {
    \node[text=ink3, anchor=west, font=\ttfamily\tiny] at (0,{1.15-\r*0.46}) {lipat \v};
    \foreach \c in {1,2,3,4} {
      \pgfmathparse{int(\c==\v)}
      \ifnum\pgfmathresult=1
        \node[va] at ({0.9+\c*1.62},{1.15-\r*0.46}) {};
      \else
        \node[tr] at ({0.9+\c*1.62},{1.15-\r*0.46}) {};
      \fi
    }
    \node[text=ink3, anchor=west] at (8.0,{1.15-\r*0.46}) {skor \v};
  }
  \node[draw=lime!50, fill=limebr!8, rounded corners=3pt, minimum width=6.5cm,
        minimum height=0.36cm, text=lime] at (4.75,-0.85) {skor akhir = rerata keempatnya};
\end{tikzpicture}

}
\end{center}
\vskip 3pt{\color{ink3}\fontsize{7.4}{9.4}\selectfont Dengan 480 sampel, satu belahan validasi tunggal terlalu berisik -- skornya bergantung pada belahan mana yang kebetulan terpilih.\par}
\icode{python}{listing 4.19 --- inti gelung K-lipat}{listings/ch04-11.py}
\ioutput{listings/ch04-12.txt}
\end{frame}

\framekicker{Bagian 4.3}
\begin{frame}[shrink=25]{Berapa lama melatihnya? Tanya kurvanya}
\icode{python}{listing 4.20 --- rerata kurva MAE}{listings/ch04-13.py}
\begin{minipage}[t]{0.485\textwidth}
\begin{iband}[signal]
MAE validasi \hilite{mendatar di sekitar epoch 120-140}, lalu mulai memburuk. Itulah titik berhentinya.
\end{iband}
\icode{python}{model akhir}{listings/ch04-14.py}

\end{minipage}%
\hfill
\begin{minipage}[t]{0.485\textwidth}
\begin{minipage}[t]{1.0\textwidth}
\istat{\textasciitilde{}0,31}{MAE uji akhir -- meleset sekitar \$31.000}
\end{minipage}%


\vskip 4pt

\begin{minipage}[t]{1.0\textwidth}
\istat{2,83}{contoh tebakan pertama $\rightarrow$ sekitar \$283.000}
\end{minipage}%

\end{minipage}%
\end{frame}

\framekicker{Ringkasan}
\begin{frame}[shrink=25]{Tabel yang akan Anda pakai terus}
\begin{itable}{>{\raggedright\arraybackslash}p{0.2632\textwidth}>{\raggedright\arraybackslash}p{0.2256\textwidth}>{\raggedright\arraybackslash}p{0.282\textwidth}>{\raggedright\arraybackslash}p{0.1692\textwidth}}
\thead{Jenis tugas} & \thead{Aktivasi lapis akhir} & \thead{Fungsi rugi} & \thead{Metrik lazim} \\ \midrule
Klasifikasi \textbf{biner} & \inlinecode{sigmoid} (1 unit) & \inlinecode{binary\_crossentropy} & accuracy \\
\textbf{Multikelas}, label one-hot & \inlinecode{softmax} (N unit) & \inlinecode{categorical\_crossentropy} & accuracy, top-K \\
\textbf{Multikelas}, label bulat & \inlinecode{softmax} (N unit) & \inlinecode{sparse\_categorical\_crossentropy} & accuracy \\
\textbf{Multilabel} & \inlinecode{sigmoid} (N unit) & \inlinecode{binary\_crossentropy} & accuracy, AUC \\
\textbf{Regresi} skalar & tanpa aktivasi (1 unit) & \inlinecode{mean\_squared\_error} & MAE \\
\end{itable}
\begin{isteps}
\item Vektorkan data diskret; \textbf{normalkan fitur dengan statistik data latih}.
\item Lapis antara \hilite{tidak boleh lebih sempit dari jumlah kelas} -- leher botol.
\item Data sedikit $\rightarrow$ model kecil, satu atau dua lapis antara.
\item Latih sampai overfit untuk \textbf{melihat} titik baliknya, lalu latih ulang sampai titik itu.
\item Data sedikit $\rightarrow$ \textbf{K-lipat}, bukan satu belahan validasi.
\end{isteps}
\vskip 4pt
\reslink{NOTEBOOK}{01\_imdb\_klasifikasi\_biner.ipynb}{../../course-slides/notebooks/ch04/01\_imdb\_klasifikasi\_biner.ipynb}
\reslink{BAB BERIKUT}{Bab 5 --- Dasar machine learning}{../ch05/index.html}
\vskip 2pt
\end{frame}

\end{document}
